\documentclass[11pt]{article}
\usepackage[T1]{fontenc}
\usepackage{lmodern}
\usepackage[margin=1in]{geometry}
\usepackage{amsmath,amssymb,amsthm,mathtools}
\usepackage{booktabs}
\usepackage{enumitem}
\usepackage{microtype}
\usepackage[colorlinks=true,linkcolor=blue,citecolor=blue,urlcolor=blue]{hyperref}

\newtheorem{theorem}{Theorem}
\newtheorem{proposition}[theorem]{Proposition}
\newtheorem{lemma}[theorem]{Lemma}
\newtheorem{corollary}[theorem]{Corollary}
\theoremstyle{definition}
\newtheorem{definition}[theorem]{Definition}
\newtheorem{remark}[theorem]{Remark}
\newtheorem{example}[theorem]{Example}

\newcommand{\C}{\mathbb{C}}
\newcommand{\R}{\mathbb{R}}
\newcommand{\F}{\mathbb{F}}
\newcommand{\inner}[2]{\langle #1,#2\rangle}
\newcommand{\norm}[1]{\lVert #1\rVert}
\newcommand{\abs}[1]{\lvert #1\rvert}
\newcommand{\op}{\mathrm{op}}
\DeclareMathOperator{\tr}{tr}
\DeclareMathOperator{\rank}{rank}

\title{\bfseries Formalising the assembly of a computer-assisted proof\\[2pt]
that no four mutually unbiased bases exist in dimension six}
\author{}
\date{14 November 2026}

\begin{document}
\maketitle

\begin{abstract}
A preprint of 24 September 2026 reports a computer-assisted proof that $\C^6$ admits no four
mutually unbiased orthonormal bases, which with the classical tensor construction gives $N(6)=3$.
Its conclusion rests on three hypotheses: a binary64 arithmetic model, a completed execution of a
100-process pipeline, and seven programs conforming to a written specification. We report on a
Lean~4 formalisation of the part of that argument that can be separated from the execution: the
reduction of the main theorem to six stage statements, and the assembly of those statements into
the final result. The assembly is proved, and --- this is the substantive finding --- its
formalisation shows that it is not the single implication the source presents but a factorisation
into two obligations: an assembly over covers, and a driver whose bookkeeping realises them. We
record what is proved, what is stated without proof, and what is not stated at all, and we separate
those three categories explicitly. Along the way we settle a disputed lemma, correct two
clause-level details of the source's statements, and report specification defects found while
transcribing them. No new axioms are assumed; every result reported is checked by the Lean kernel,
and we list the 23 remaining holes rather than summarising them away.
\end{abstract}

\tableofcontents

\section{Introduction}

\subsection{The result and its trust structure}

Let $N(d)$ be the largest size of a family of pairwise mutually unbiased orthonormal bases of
$\C^d$. The tensor construction gives $N(6)\ge 3$, and a preprint of 24 September 2026
\cite{source} reports the matching upper bound. The claim is computer-assisted, and the source is
explicit that it is conditional. Reduced to essentials, the conclusion rests on:
\begin{enumerate}[label=(\Alph*)]
\item a model of binary64 arithmetic, of the compiler, and of the C library's elementary functions;
\item a completed execution of a pipeline of roughly one hundred processes, whose records are
      archived and hash-identified;
\item seven programs that conform to the specification written in the paper's Sections~2--6.
\end{enumerate}
None of these is a mathematical hypothesis, and a reader cannot check any of them by reading. That
is the situation our formalisation addresses, on the part of the argument where it can be addressed:
the \emph{assembly}, that is, the step from local statements about finitely many regions to the
global conclusion.

\subsection{What we formalised, and what we deliberately did not}

The source's argument has two layers. The \emph{analytic} layer consists of necessary conditions
and covering estimates: each stage produces a finite family of regions that covers the objects a
counterexample would need, and each region carries a certificate that excludes it or contains it.
The \emph{assembly} layer is the composition of those coverings. We formalised the assembly in full
and the analytic layer only in the parts that are elementary and self-contained; the analytic middle
(notably the refinement relations of the source's Section~5) we did not attempt, and we say so where
it matters rather than implying completeness.

The distinction is not a convenience. The assembly is where a wrong proof of this shape fails
silently: an error in the induction over regions, in the quantifier order of a covering statement,
or in the handling of the equivalence relation under which regions are compared, produces a proof
that reads correctly and establishes nothing. Those are precisely the errors a machine catches.

\subsection{Contributions}

\begin{enumerate}[label=\textbf{C\arabic*.}]
\item A formal statement of the source's main theorem as an \emph{equivalence} with six stage
      statements together with seven terminal predicates, rather than as an implication
      (Section~\ref{sec:assembly}). The equivalence is what makes the obligations visible.
\item A proof of the assembly, and of the two inductions underneath it --- over subdivision depth,
      and over the order of completed parents --- in a form that cannot assume what it proves
      (Section~\ref{sec:assembly}).
\item The finding that the source's Theorem~7.1 requires a \emph{driver} hypothesis it does not
      state: the seven coverage facts are not consequences of the paper's bookkeeping conditions,
      so the theorem is two steps where the paper presents one (Section~\ref{sec:findings}).
\item Formal proofs of the fine-grained lemmas that the rest of the argument consumes, including
      the double-centring projection in full, the stability of circular gaps, the operator
      estimate behind the refinement stage, and the Gauss sums underlying the lower bound
      (Section~\ref{sec:proved}).
\item An adjudication of a lemma whose status was in dispute, with a counterexample that turns out
      to violate its own hypothesis (Section~\ref{sec:l51}).
\item A status table that separates proved, stated-without-proof, and not-stated
      (Section~\ref{sec:status}), and a list of the specification defects found while preparing the
      formal statements (Section~\ref{sec:spec}).
\end{enumerate}

\subsection{Relation to other work}

An independent proof of the same theorem appeared on 4 November 2026 \cite{WT26}. It proceeds by a
different route --- a connectedness theorem for order-six complex Hadamard matrices, followed by
interval-arithmetic certification over overlapping phase balls --- and it also classifies all sets
of mutually unbiased bases in dimension six. The two proofs have different trust bases: ours and the
source's rest on a binary64 model and an execution, that one on directed-rounding interval
arithmetic. Neither subsumes the other, and this paper takes no position on which should be
preferred; we mention it because any statement of the source's status that omits it is now
incomplete.

\section{The source argument in brief}
\label{sec:source}

Fix a pair of mutually unbiased bases and send the first to the coordinate basis by a common
unitary. The second is then $H/\sqrt 6$ for a complex Hadamard matrix $H$, that is, a matrix with
entries of modulus one and $HH^*=6I$. The upper bound is the statement that no such $H$ extends to
two further bases. The source's argument covers the second basis by finitely many balls and excludes
extension uniformly on each:

\begin{description}[leftmargin=1.2em,style=nextline]
\item[Seed cover.] Row and column rephasings reduce $H$ to nine free phases, chosen through a
  minimum circular gap among the phase-difference multisets of all row pairs and column pairs. Those
  nine phases are covered by finitely many closed balls of radius $0.032$.
\item[Completion.] Splitting $H$ into blocks of sizes $2+4$, the fixed seed determines two rank-two
  kernels; compressing the unknown lower-right block between them yields a $2\times2$ matrix whose
  polar factor is a unitary. A closed cover of $U(2)$ therefore accounts for every feasible
  completion, and Taylor estimates turn each cell into either an exclusion or a certified ball.
\item[Refinement.] Each certified ball is refined, using tangent coordinates and a certified
  estimate of the normal displacement, into balls of a prescribed smaller radius.
\item[Exclusion.] Inside a final ball, all vectors unbiased to the first two bases are covered by
  small balls; two graphs record which pairs of balls can be orthogonal and which can be unbiased;
  two further bases would require a specific pattern of two disjoint six-cliques with all cross
  edges present, and a finite search excludes that pattern.
\end{description}

The assembly, the source's Theorem~7.1, is then: if the pipeline's terminal conditions hold, and
each stage delivers its covering, no four bases exist. Our contribution concerns that last sentence.

\begin{remark}
The source's Lemma~5.1 is used at the hinge of the refinement stage, and its status was disputed
within our project; Section~\ref{sec:l51} settles it. Its Lemma~3.1 carries a hypothesis
$\abs S<4$ that a reader may take for a technicality; it is not (Section~\ref{sec:findings}).
\end{remark}

\section{The formal development}

\subsection{Setup}

The development is written in Lean~4 (version 4.33.1) against Mathlib. It introduces no axioms: every
declaration reported as proved depends only on \texttt{propext}, \texttt{Classical.choice} and
\texttt{Quot.sound}, and every hole is a \texttt{sorry} that we list explicitly, so that
\texttt{sorryAx} appears in a declaration's axiom footprint exactly when that declaration is
incomplete. The development is 16 modules with an acyclic import graph; reusable material (finite
circular multisets, the double-centring projection, the Hadamard predicate, operator bounds for
nearly-orthogonal frames) is separated from paper-specific assembly.

Two conventions were fixed early, both of which matter for fidelity.

\begin{definition}[Hadamard]
A matrix $H\in \C^{6\times 6}$ is \emph{Hadamard} if every entry has squared modulus one and
$HH^*=6I$. The condition $H^*H=6I$ is \emph{not} part of the definition.
\end{definition}

This is the source's Definition~2.1, and it is deliberately one-sided. For a square matrix the
two-sided condition follows, but Mathlib's \texttt{Matrix.IsHadamard} is two-sided, and importing it
would silently strengthen the hypothesis. A formalisation that did so would prove a true theorem
about a different class of objects.

The second convention concerns the equivalence under which regions are compared. Throughout, a
stage does not cover a set of matrices but a set of \emph{equivalence classes} of matrices, the
equivalence being generated by row and column permutations, rephasings, conjugation and
transposition. We therefore model membership as membership of a class, not of a matrix. This is not
a stylistic choice: it is what allows the ``up to equivalence'' clause of every stage to compose,
and a formalisation that used set membership would need an additional compatibility lemma at each
of the six junctions.

\subsection{The interface}
\label{sec:assembly}

The assembly is stated over an abstract type of regions. Let $\mathtt{Pt}$ be a type of points,
$\mathtt{emb}$ a relation on it intended as equivalence of phase matrices, \texttt{ext} a predicate
intended as extendibility by two further bases, and \texttt{Cov emb A B} the assertion that every
point of $A$ is equivalent to a point of $B$. Each stage statement then has the same shape:
\begin{center}
\texttt{Cov emb (input set) (output set)}.
\end{center}
The two inductions that the source performs in prose are captured by an inductive certificate type,
which is what prevents circularity:
\begin{description}[leftmargin=1.4em,style=nextline]
\item[\texttt{empty}] the target and source sets are disjoint;
\item[\texttt{excl}] a certificate that the region excludes every point of $S$;
\item[\texttt{ball}] a certified ball $b$ of the region family with $S\subseteq\mathtt{mem}(b)$;
\item[\texttt{union}] an exhaustive decomposition of $S$ into finitely many sub-regions.
\end{description}
and the elimination principle \texttt{Cert.covered} gives $\mathtt{Covered}\,(T\cap S)$ from a
certificate. The blocker induction --- the step in which a completed parent is used to certify a
later region, which is the one place where a proof of this shape can assume its own conclusion ---
is a separate theorem \texttt{accumulate\_covered}, whose hypothesis is a certificate at the prefix
$\ell \mathbin{+\!+} [p]$ of the region list rather than at the whole list.

\begin{theorem}[assembly]\label{thm:assembly}
Let $H$ be the set of Hadamard phase matrices, \texttt{ext} be invariant under $\mathtt{emb}$, and
let $S_{\mathrm{seed}},S_{\mathrm{coarse}},S_{\mathrm{mid}},S_{\mathrm{fine}},
S_{\mathrm{init}},S_{\mathrm{fallback}},S_{\mathrm{final}}$ be sets of regions. If
\begin{align*}
&\mathtt{Cov}\,H\,S_{\mathrm{seed}} &&\text{(seed cover)}\\
&\mathtt{Cov}\,S_{\mathrm{seed}}\,S_{\mathrm{coarse}} &&\text{(completion)}\\
&\mathtt{Cov}\,S_{\mathrm{coarse}}\,S_{\mathrm{mid}},\quad
 \mathtt{Cov}\,S_{\mathrm{mid}}\,S_{\mathrm{fine}} &&\text{(refinement)}\\
&\mathtt{Cov}\,S_{\mathrm{fine}}\,S_{\mathrm{init}},\quad
 \mathtt{Cov}\,S_{\mathrm{init}}\,S_{\mathrm{fallback}},\quad
 \mathtt{Cov}\,S_{\mathrm{fallback}}\,S_{\mathrm{final}} &&\text{(exclusion)}\\
&\textstyle\bigwedge_{x\in S_{\mathrm{final}}}\neg\,\mathtt{ext}(x) &&\text{(no extension)}\\
\end{align*}
then $\neg\,\mathtt{ext}(x)$ for every $x\in H$.
\end{theorem}

The proof is transitivity of \texttt{Cov} composed seven times, followed by
\texttt{Cov.excludes}. It is short, and that is the point: all the difficulty is in the seven
hypotheses, and the theorem's formal statement forces them to be named.

\begin{proposition}[the assembly is an equivalence]\label{prop:iff}
The source's main theorem is equivalent to the seven covering statements of
Theorem~\ref{thm:assembly} together with the seven terminal predicates of its Section~7.2.
\end{proposition}

Proposition~\ref{prop:iff} is the formal counterpart of the observation that motivates the paper.
The source presents one theorem; there are in fact two kinds of obligation in it, and separating
them is what makes the second one (Section~\ref{sec:findings}) visible.

\section{What formalisation found}
\label{sec:findings}

\subsection{Theorem 7.1 needs a driver}

The hypothesis list of Theorem~\ref{thm:assembly} is exactly the seven coverage facts and the
terminal exclusion. It contains \emph{no} bookkeeping condition. But the source's Theorem~7.1 is
stated under its Section~7.2, which requires builds to succeed, specified exit codes, empty failure
files, batch counts to agree, and the fallback partition to be exhaustive. Those conditions are
neither hypotheses nor consequences of the coverage facts: they are what makes the coverage facts
\emph{true of the records the run produced}.

The formal version therefore factors the source's single theorem into
\begin{enumerate}[label=(\roman*)]
\item the assembly (Theorem~\ref{thm:assembly}), which is mathematics; and
\item a driver lemma --- call it \emph{bookkeeping identifies records} --- which says that a
      pipeline whose control flow and record counts satisfy Section~7.2 produces a region list on
      which each of the seven covering statements holds.
\end{enumerate}
The source's proof of Theorem~7.1 implicitly uses (ii) and never states it. This is the most
consequential finding of the formalisation, because (ii) is exactly the step where a silently
dropped region would invalidate the conclusion, and it is the step that a checker for the
certificate would have to implement.

\subsection{Two clause-level corrections}

\textbf{The `up to equivalence' clause must be class membership.} Six of the seven covering
statements in the source are qualified by ``up to the equivalences of Lemma~2.2''. If these are read
as set inclusions, the composition fails: the output of one stage is not a set of matrices but a set
of classes. Read as class membership, they compose, and the extendibility-invariance of the
equivalence (a separate lemma, proved) supplies the only compatibility that is needed. The source
never states the joint lemma; its proof of Theorem~7.1 uses it at every junction.

\textbf{Both cliques in the exclusion pattern are orthogonality cliques.} The configuration that
two further bases would force is two disjoint sets $A,B$ of ball labels, each of six labels, each a
clique in the orthogonality graph, with every cross pair joined in the unbiasedness graph. In the
source this is one displayed formula, and it is easy to read the second condition as a
unbiasedness clique on $B$; it is not, and the search is stated for the correct condition. During
formalisation we wrote it the wrong way first, which is evidence of how much work that single
formula is doing.

\subsection{The hypothesis of Lemma 3.1 is load-bearing}

The source's Lemma~3.1 states that the squared singular values of $M=(1\ \ p)$ are $4\pm\abs S$,
where $S=\sum_j p_j$, \emph{if} $\abs S<4$, and constructs a $4\times2$ orthonormal frame $K$ whose
range is $\ker M^*$. At $\abs S=4$ one has $p=3\cdot\mathbf 1$, so $M$ has rank one, the singular
values are $8$ and $0$, and $\ker M^*$ is three-dimensional --- no $4\times 2$ frame can have that
range. The hypothesis is therefore not a convenience but the condition under which the dimension
count works, and the completion stage's frame construction fails silently without it.

\subsection{Definition 2.1 must stay one-sided}
Formalising with a two-sided Hadamard predicate would be unsound bookkeeping even though the two
conditions are equivalent for square matrices: the source's hypothesis is $HH^*=6I$, and the
equivalence itself is a (small) theorem that should not be assumed.

\section{The disputed lemma}
\label{sec:l51}

The refinement stage rests on the following.

\begin{lemma}[source, Lemma 5.1]\label{lem:51}
Let $T:\R^k\to\R^{36}$ satisfy $\norm{T^*T-I}\le o$, put $A=I-TT^*$, and for $x\in\R^{36}$ write
$\zeta=T^*x$. Then $\norm T\le\sqrt{1+o}$, $\norm A\le 1+o$, and
\begin{equation}\label{eq:51}
\norm{x}^2\le (1+o)\norm{\zeta}^2+\norm{Ax}^2 .
\end{equation}
No horizontality of $x$ or of the columns of $T$ is required.
\end{lemma}

Within the project this lemma was reported to be false, with the witness $T=\varepsilon P$ for a
rank-one projection $P$, $o=0$, and $x=e$. The witness violates the hypothesis: $o=0$ forces
$\norm{T^*T-I}=0$, hence $\varepsilon=1$, and then \eqref{eq:51} holds with equality. The lemma is
true, and here is the derivation:
\begin{align*}
\norm{x-T\zeta}^2
&=\norm x^2-2\inner{x}{T\zeta}+\norm{T\zeta}^2
 =\norm x^2-2\norm\zeta^2+\norm{T\zeta}^2,
\end{align*}
using $\inner{x}{T\zeta}=\inner{T^*x}{\zeta}=\norm\zeta^2$, and $\norm{T\zeta}^2\ge(1-o)\norm\zeta^2$
by the hypothesis. Hence
$\norm{x-T\zeta}^2\ge\norm x^2-(1+o)\norm\zeta^2$, which is \eqref{eq:51} since $x-T\zeta=Ax$.
The formal development proves both the two-sided Gram consequence
$(\forall y)(1-o)\norm y^2\le\norm{Ty}^2$ and \eqref{eq:51} from the hypothesis alone, with no
additional binder. Two lessons: a near-isometry hypothesis is two-sided, and a counterexample to a
lemma must be checked against the lemma's hypothesis before it is reported.

\section{What is proved}
\label{sec:proved}

We distinguish three categories, and keep them separate on purpose: a statement is \emph{proved}
(kernel-checked, no holes), \emph{stated} (transcribed faithfully, with a named hole), or
\emph{absent} (not attempted). Conflating the last two is how a formalisation report overstates.

\begin{description}[leftmargin=1.2em,style=nextline]
\item[Proved, in full.] The assembly of Theorem~\ref{thm:assembly} and Proposition~\ref{prop:iff};
  the certificate algebra and \texttt{accumulate\_covered}; the double-centring projection
  $\Pi$ --- idempotence, range equal to the horizontal subspace, the decomposition of $z-\Pi z$
  into a row-constant plus a column-constant part with explicit witnesses, the Frobenius
  contraction, and the packaged rephasing $e(P+\Pi(Q-P))$ as a product of row and column phases;
  the stability of circular gaps, in three forms (arc transfer, gap-bound transfer, and the literal
  $\pm2\delta$ statement); Lemma~\ref{lem:51} and its Gram consequence; the Gauss sums underlying
  the lower bound, namely $\sum_{x<3}\omega^{x^2+tx}\in\{1+2\omega,\,2+\omega^2\}$ with
  $\omega=e^{2\pi i/3}$, each of squared modulus three; the exclusion-side lemmas on distinct labels
  and on the graph pattern; and the norm-exclusion inequality of the completion stage.
\item[Stated, with named holes.] The six clauses of the source's Lemma~2.2 and their packaged form;
  the frame lemma beyond its operator statement; the directional-exclusion inequality of the
  completion stage; the ellipsoid containment and offset bounds in their remaining goals; the
  identity that the cross-product entries of the lower bound have squared modulus $1/3$; and the
  existence of three mutually unbiased bases.
\item[Absent.] The source's Section~5 in its entirety --- the certified relations between two
  centres, the distance iteration, and the tangent grid. This is the largest gap in the development
  and is a deliberate consequence of the scope we set: the analytic middle is where the floating
  point enters, and formalising it without a decision about the arithmetic model would have
  produced statements that could not be discharged.
\end{description}

\begin{remark}
The existence of three bases is the one part of the main theorem that is not computer-assisted, and
it is the part a reader can check by hand. It is worth finishing in the formal development precisely
because it is unconditional: it would give a kernel-checked $N(6)\ge 3$ next to a conditional
$N(6)\le 3$, and the asymmetry between the two is the paper's real content.
\end{remark}

\section{Status of the development}
\label{sec:status}

Table~\ref{tab:status} records the hole inventory. Sixteen modules elaborate; twenty-three
declarations carry a hole, distributed over seven of them. The nine modules that carry none fall
into three groups that must not be confused. \texttt{Interface} and \texttt{AssemblyCover} are
hole-free because the assembly is complete. \texttt{Basic}, \texttt{Gap}, \texttt{Hypotheses} and
\texttt{AxiomCheck} are hole-free because their content is finished and elementary. And \texttt{Completion}, \texttt{Refine} and \texttt{Assembly} are \emph{stubs} that
assert nothing at all. A reader scanning a hole count will read those zeros as finished stages;
they are not, and we list them separately for that reason.

\begin{table}[ht]
\centering
\small
\begin{tabular}{@{}llp{7.4cm}@{}}
\toprule
module & holes & content \\
\midrule
\texttt{Basic} & 0 & horizontal projection $\Pi$, rephasing, integer shifts \\
\texttt{Hadamard} & 7 & Lemma 2.2's six clauses and their packaged form \\
\texttt{Seed} & 1 & stability of circular gaps (a second copy; see below) \\
\texttt{ThreeMUB} & 2 & Gauss sums; existence of three bases \\
\texttt{Frames} & 7 & Lemma 3.1 beyond the operator statement \\
\texttt{Exclusion} & 1 & directional exclusion (the source's \texttt{eq} layout) \\
\texttt{Ellipsoid} & 4 & offset ellipsoid bound, remaining goals \\
\texttt{Exclude} & 1 & bin displacement of the vector cover (Lemma 6.1) \\
\texttt{Interface}, \texttt{AssemblyCover} & 0 & assembly: \emph{proved} \\
\texttt{Completion}, \texttt{Refine}, \texttt{Assembly} & 0 & \emph{stubs: assert nothing} \\
\midrule
total & 23 & of which 3 modules are stubs \\
\bottomrule
\end{tabular}
\caption{Declared holes. A hole is a \texttt{sorry}; \texttt{sorryAx} then appears in the axiom
footprint of exactly the incomplete declarations.}
\label{tab:status}
\end{table}

Two bookkeeping defects in the development itself are worth recording, since they affect how the
table should be read. Lemma~2.5 is proved in one module and simultaneously present as a hole in
another, so the count is pessimistic where the proof already exists. And the codebase contains
three stub modules that elaborate with no holes for the trivial reason that they state nothing;
they are placeholders for the refinement stage, not evidence about it.

\section{Specification defects found while transcribing}
\label{sec:spec}

Preparing faithful formal statements requires reading the specification more closely than a reader
does, and this produced findings about the source's parameters and prose. None of them affects the
truth of the source's conclusion; all of them affect what the paper should say. We report them
because a source that is to be verified should be unambiguous, and because several are the kind of
statement a reader would take on trust.

\textbf{The bin count.} The source fixes $N=112$ in one sentence and offers one reason: that
$\binom{116}{5}$ is below the signed 32-bit limit. That reason is inert --- the same paragraph says
the traversal uses 64-bit counters --- and it is not binding: the largest $N$ it permits is 189. The
constraints the source's own mechanism imposes are sharper. Each uncovered bin must fit in a ball
of radius $R=0.032$, giving $N\ge 47$. The five-coordinate bin differences are stored in a signed
eight-bit type, giving $N\le 128$. The reuse test can be guaranteed to succeed only when
$\norm k\le NR-\tfrac32$, which first admits two-step offsets at $N\ge110$; $112$ is the smallest
multiple of $16$ above that threshold. And the reuse threshold is written with $112$ hardwired where
$N$ belongs, so the pipeline \emph{as printed} is sound only for $N\ge112$. The value is therefore
defensible --- it is the left endpoint of the window it is itself consistent with --- but not for
the reason printed.

\textbf{The refinement widths.} The base grid of the refinement stage has $(2L+1)^k$ cells with
$L=\lceil E/w+\tfrac12\rceil$, capped at $5\cdot10^7$. At the three specified widths the admissible
root dimension is $k\le5,5,7$, not the stated $k\le10$; a record with $k=6$ entering the second
pass cannot be refined at all. Since the run reports no failed parent, every record must satisfy
$k\le5$, and the stated allowance is half unused.

\textbf{A clustering cutoff that does no work.} The vector stage's clustering accepts a candidate
centre when its effective radius reaches the spacing of bin midpoints, which requires $M\ge51$; the
three coarser grids are $12,24,48$. On those grids the cutoff cannot admit any merge, whatever the
sign convention, and the clustering stage is a no-op.

\textbf{A sentinel inside its own range.} One routine signals failure by returning $100$, while the
same appendix proves the quantity being bounded is at most $100.1$. The sentinel is therefore not
distinguishable from a legitimate value, and safety comes from a later unit conversion rather than
from the value itself.

\textbf{An undefined count.} Table~1 of the source reports $151{,}744{,}086$ ``pairs''. Three
independent replays agree that no reading of the stated filters produces that number: the
five-coordinate enumeration admits $85{,}349{,}292$ lists, and after the shared-displacement
intersection the tail side admits at most $112{,}160{,}350$ pairs. One replay reproduces the
source's retained count of $22{,}829$ lists exactly, which localises the discrepancy to the meaning
of the word ``pairs'' rather than to any filter. The paper should say which product it is counting.

\section{Discussion}

\subsection{What this establishes, and what it does not}

We have formalised the assembly of the source's argument and the elementary lemmas it consumes. We
have \emph{not} formalised the analytic middle, and we have \emph{not} verified the source's
execution. Consequently this paper does not, by itself, establish that $N(6)=3$; it establishes that
the source's conclusion follows from seven explicitly stated covering facts plus a driver lemma,
that the source's own presentation of that step is missing one of the two, and that the lemmas on
which the covering facts rest are correct where we formalised them.

That is a real gain in a specific sense: the residual trust is now located. After this work, a
reader who doubts the conclusion can direct attention to (i) the driver lemma, which is finite
combinatorics over records and is the natural target for a checker, and (ii) the analytic middle,
which is where the arithmetic model enters. Before this work the two were fused into a single
computer-assisted claim.

\subsection{The arithmetic question}

The source converts a binary64 execution into a theorem by proving one-sided error bounds for every
certificate evaluation, roughly twenty pages of them. Those bounds are correct as far as we can
check, and they are fragile in a specific way: they are tied to an evaluation order and to a
compiler, and one of them required inspecting a disassembly to confirm that two power expressions
were lowered to multiplications. We did not attempt to formalise them. The alternative --- evaluating
certificates in directed-rounding interval arithmetic, so that ``the stored number is an upper
bound'' becomes a theorem of a library rather than of a page of case analysis --- would reduce the
arithmetic appendix to a few pages and remove the compiler dependency. The independent proof
\cite{WT26} takes exactly that route, which suggests the replacement is practical.

\subsection{Threats to the fidelity of this report}

The statements we formalised were transcribed by hand from the source's text, and the report's
fidelity claim is therefore limited by that transcription. In the course of the work, four
formula-level claims about the source were circulated within the project and then withdrawn: all
four concerned radicals and bracket extents, and all four were artefacts of text extraction rather
than defects in the source. The lesson is general and worth stating: no claim about a displayed
formula should be made from an extracted text; radicals, overlines and superscript placement must be
checked on a rendered page. Every formula-level claim in this paper has been checked that way.

\section{Conclusion}

The assembly of the source's argument is formalised, and its formalisation shows that the assembly
is not the single step the source presents: it is an assembly over covers plus a driver that
realises the covers from a run, and only the first is mathematics. The elementary lemmas are
correct where formalised; one disputed lemma is settled in the source's favour; two clause-level
details need correction; and five specification defects are recorded. The analytic middle and the
execution remain unverified here, and the honest summary of the source's status is that it is
conditional on three hypotheses, one of which --- the arithmetic model --- has an alternative that
has since been implemented by an independent proof of the same theorem.

\begin{thebibliography}{99}
\bibitem{source} OpenAI. \emph{The maximum number of mutually unbiased bases in dimension six.}
  Preprint, 24 September 2026.
\bibitem{WT26} M.~C\'ardenes Wuttig and J.~Tindall. \emph{Using topology to reveal all mutually
  unbiased bases in dimension six.} arXiv:2611.06917 [quant-ph], 4 November 2026.
\bibitem{Zauner} G.~Zauner. \emph{Quantendesigns: Grundz\"uge einer nichtkommutativen
  Designtheorie.} PhD thesis, Universit\"at Wien, 1999.
\bibitem{KR} A.~Klappenecker and M.~R\"otteler. \emph{Constructions of mutually unbiased bases.}
  Finite Fields and Applications, LNCS 2948, 2004.
\bibitem{JMMSW} P.~Jaming, M.~Matolcsi, P.~M\'ora, F.~Sz\"oll\H{o}si and M.~Weiner.
  \emph{A generalized Pauli problem and an infinite family of MUB-triplets in dimension 6.}
  J. Phys. A 42 (2009).
\bibitem{Grassl} M.~Grassl. \emph{On SIC-POVMs and MUBs in dimension 6.} arXiv:quant-ph/0406175,
  2004.
\bibitem{BH} P.~Butterley and W.~Hall. \emph{Numerical evidence for the maximum number of mutually
  unbiased bases in dimension six.} Phys. Lett. A 369 (2007).
\bibitem{BW} S.~Brierley and S.~Weigert. \emph{Constructing mutually unbiased bases in dimension
  six.} Phys. Rev. A 79 (2009).
\bibitem{Bengtsson} I.~Bengtsson et al. \emph{Mutually unbiased bases and Hadamard matrices of order
  six.} J. Math. Phys. 48 (2007).
\bibitem{BSTW} P.~O.~Boykin, M.~Sitharam, P.~H.~Tiep and P.~Wocjan. \emph{Mutually unbiased bases
  and orthogonal decompositions of Lie algebras.} arXiv:quant-ph/0506089, 2005.
\bibitem{Lean} L.~de Moura and S.~Ullrich. \emph{The Lean 4 theorem prover and programming
  language.} CADE-28, 2021.
\bibitem{Mathlib} The Mathlib Community. \emph{The Lean mathematical library.} CPP 2020.
\end{thebibliography}

\appendix

\section{Selected Lean signatures}
\label{app:sig}

The following are transcribed from the development, with notation lightened and hypotheses
elided where they are not the point.

\begin{quote}\ttfamily\small
theorem Cov.trans : Cov emb A B $\to$ Cov emb B C $\to$ Cov emb A C\\[3pt]
theorem Cert.covered : Cert mem B T S $\to$ Covered (T $\cap$ S)\\[3pt]
theorem accumulate\_covered :\\
\ \ (h : $\forall$ p $\in$ ps, Cert mem B T (mem p)) $\to$ Covered (T $\cap$ $\bigcup$ ps)\\[3pt]
theorem theorem\_7\_1 (emb) (htrans) (ext) (hinv)\\
\ \ (H Sseed Scoarse Smid Sfine Sinit Sfallback Sfinal)\\
\ \ (h2\_10 : Cov emb H Sseed) (hA : Cov emb Sseed Scoarse)\\
\ \ (hBmid : Cov emb Scoarse Smid) (hBfine : Cov emb Smid Sfine)\\
\ \ (hinitC : Cov emb Sfine Sinit) (hfallback : Cov emb Sinit Sfallback)\\
\ \ (hfinal : Cov emb Sfallback Sfinal)\\
\ \ (hCneg : $\forall$ y $\in$ Sfinal, $\neg$ ext y) :\\
\ \ $\forall$ x $\in$ H, $\neg$ ext x\\[3pt]
theorem lemma51\_third (T) (o) (h : GramNear T o) (x) :\\
\ \ $\norm{x}^2 \le (1+o) * \norm{T^* x}^2 + \norm{x - T(T^* x)}^2$
\end{quote}

\begin{remark}
The signature of \texttt{theorem\_7\_1} is the evidence for the paper's main finding: it takes
exactly the seven coverage hypotheses and the terminal exclusion, and nothing else. The source's
Section~7.2 conditions appear nowhere in it.
\end{remark}

\end{document}
